% Version 3.1 December 2024
% Based on the official Springer Nature LaTeX authoring template.
% Draft status: internal BioCC manuscript scaffold, not submission-ready.

\documentclass[pdflatex,sn-nature]{sn-jnl}

\usepackage{graphicx}
\usepackage{multirow}
\usepackage{amsmath,amssymb,amsfonts}
\usepackage{booktabs}
\usepackage{array}

\raggedbottom

\begin{document}

\title[Interpretable protein-model perturbation diagnostics]{Interpretable protein language model features for variant perturbation and generative circuit diagnostics}

\author*[1]{\fnm{Zipeng} \sur{Liu}}\email{liuzipeng@tju.edu.cn}
\author[1]{\fnm{Coauthors} \sur{to be added}}

\affil*[1]{\orgdiv{Affiliation to be updated}, \orgname{Institution to be updated}, \orgaddress{\city{City}, \country{Country}}}

\abstract{Protein language models provide accurate scalar variant scores, but scalar scores often do not explain which molecular function is perturbed or whether a proposed intervention changes a generative mechanism. We developed BioCC, a sparse-feature analysis workflow for protein language models that converts hidden activations into feature-level perturbation signatures, annotation-linked residue examples and circuit-level diagnostics. In ESM-2 variant analysis, annotation-selected sparse autoencoder features complemented masked-marginal log-likelihood for ClinVar and cancer pathogenicity prediction, but did not exceed AlphaMissense and did not generalize as a robust LOF/GOF/DN classifier under protein-level holdout. In generative protein models, circuit-latent dictionaries exposed enzyme-class features and enabled controlled-generation tests, but current steering effects were not statistically significant. These results position BioCC as an interpretable diagnostic and hypothesis-generation framework rather than a finished universal variant-effect or protein-design predictor.}

\keywords{protein language models, sparse autoencoders, variant interpretation, missense variants, indels, circuit tracing, protein generation}

\maketitle

\section{Introduction}\label{sec:introduction}

Protein language models (PLMs) have become useful zero-shot models of protein sequence constraint and variant effect \cite{lin2023esm,meier2021esm1v}. Their masked-marginal or ensemble likelihood scores can rank deleterious substitutions without supervised training, and specialized models such as AlphaMissense now provide strong proteome-wide pathogenicity estimates for human missense variants \cite{cheng2023alphamissense}. These scalar predictors are clinically useful, but they leave two gaps. First, a scalar score does not usually explain whether a variant perturbs folding, an active site, a binding interface, regulation or another functional axis. Second, in generative PLMs, a high-level class prompt or steering intervention may change sampled sequences without proving that a causal biological feature has been controlled.

Sparse autoencoders (SAEs) and circuit-latent dictionaries offer a practical way to inspect intermediate model representations. They decompose dense activations into sparse feature activations that can be mapped to top-firing proteins, residue positions and annotation categories. In principle, the difference between wild-type and mutant sparse activations can describe a variant as a vector-valued perturbation rather than as one score. Similarly, sparse features in generative models can be tested by activation interventions during sequence generation. The key question is whether these interpretable signatures add experimentally useful information beyond existing scalar predictors.

Here we report an initial, deliberately conservative evaluation of BioCC, a two-part workflow for interpretable protein-model analysis. The first part uses ESM-2-3B SAEs across layers 19, 23, 27, 31 and 35 to score ClinVar, cancer holdout, ProteinGym and indel variants. The second part trains circuit-latent dictionaries for ProtGPT2 and ZymCTRL and evaluates enzyme-class steering, structural quality and cross-model feature conservation. We emphasize both positive and negative results. BioCC improves the local ESM-2 LLR pathogenicity baseline when combined with annotation-selected SAE features, and it produces feature-level perturbation tables that can guide experimental follow-up. However, AlphaMissense remains stronger for scalar pathogenicity, protein-level mechanism classification is near random, ProteinGym does not show a mean advantage over LLR, and current generative steering experiments are negative. These outcomes define the current value proposition: BioCC is a mechanism-like diagnostic and wet-lab prioritization tool, not yet a validated mechanism classifier or design engine.

\section{Results}\label{sec:results}

\subsection{Sparse protein features align with residue-level annotations}

We trained ESM-2-3B sparse autoencoders on five production layers and linked feature firing positions to Swiss-Prot, Pfam, Gene Ontology and binding-site annotations. The firing-position rerun increased the number of annotation-aligned features at all layers. At a feature-level F1 threshold of 0.5, known-feature counts increased from 146 to 381 at layer 19, from 135 to 301 at layer 23, from 72 to 163 at layer 27, from 49 to 86 at layer 31 and from 32 to 45 at layer 35 (Table~\ref{tab:annotation}). The deepest layer improved but did not meet the internal target of at least 60 known features.

\begin{table}[h]
\caption{Annotation alignment after the firing-position rerun. Known features use F1 $>$ 0.5, partial features use F1 $>$ 0.2 and useful features use F1 $>$ 0.1.}\label{tab:annotation}
\begin{tabular}{@{}rrrrrrr@{}}
\toprule
Layer & Known before & Known after & Partial after & Useful after & Features with firing positions \\
\midrule
19 & 146 & 381 & 3483 & 7481 & 13868 \\
23 & 135 & 301 & 2992 & 6773 & 13832 \\
27 & 72  & 163 & 1725 & 4423 & 13456 \\
31 & 49  & 86  & 1135 & 3197 & 13649 \\
35 & 32  & 45  & 653  & 1935 & 12895 \\
\botrule
\end{tabular}
\end{table}

The feature audit produced 150 top coefficient-linked rows across DN, GOF and LOF labels, five layers and ten top features per class-layer pair. Each row records the logistic regression coefficient, feature kind, best annotation, expanded Pfam or functional proxy and top residue-level firing examples. This table is usable for figure planning and manual residue triage. It is not direct biological validation: many features remain family-level or weakly annotated, so mechanistic language should be framed as a hypothesis to be tested rather than as a settled assignment.

\subsection{SAE perturbations complement ESM-2 LLR but do not beat AlphaMissense}

On 2,000 ClinVar missense variants, annotation-selected SAE features alone reached AUC 0.8782 for pathogenicity, close to the ESM-2 masked-marginal LLR baseline of 0.8822. A simple z-sum ensemble of SAE-LR and LLR improved to AUC 0.9143. On a 101-variant cancer holdout, SAE-LR reached AUC 0.9079, LLR reached 0.8978 and the ensemble reached 0.9193 (Table~\ref{tab:pathogenicity}). These results show that SAE perturbation signatures carry information not fully captured by LLR.

\begin{table}[h]
\caption{Pathogenicity performance on local R1 cohorts. Confidence intervals are bootstrap 95\% intervals.}\label{tab:pathogenicity}
\begin{tabular}{@{}llllr@{}}
\toprule
Cohort & Method & AUC & 95\% CI & n \\
\midrule
ClinVar2000 & SAE-LR & 0.8782 & [0.8619, 0.8923] & 2000 \\
ClinVar2000 & ESM-2 LLR & 0.8822 & [0.8677, 0.8960] & 2000 \\
ClinVar2000 & SAE+LLR & 0.9143 & [0.9010, 0.9259] & 2000 \\
ClinVar2000 & AlphaMissense & 0.9474 & [0.9377, 0.9567] & 1972 \\
CancerHoldout101 & SAE-LR & 0.9079 & [0.8471, 0.9590] & 101 \\
CancerHoldout101 & ESM-2 LLR & 0.8978 & [0.8273, 0.9561] & 101 \\
CancerHoldout101 & SAE+LLR & 0.9193 & [0.8617, 0.9646] & 101 \\
CancerHoldout101 & AlphaMissense & 0.9700 & [0.9244, 0.9988] & 101 \\
\botrule
\end{tabular}
\end{table}

However, the same comparison also narrows the claim. AlphaMissense matched 1,972 of 2,000 ClinVar variants and all 101 cancer-holdout variants, with AUC 0.9474 and 0.9700, respectively. BioCC therefore should not be presented as a state-of-the-art scalar pathogenicity predictor on these cohorts. Its remaining novelty must come from interpretable perturbation structure, residual analysis against scalar tools and variant classes where missense-only tools are less uniformly applicable.

\subsection{Mechanism classification is inflated by variant-level cross-validation}

We built a mechanism-labeled subset of 497 ClinVar variants covering 80 DN, 120 GOF and 297 LOF labels across 255 proteins. Under variant-level cross-validation, SAE-only features achieved macro-AUC 0.7471, while LLR alone was near random at 0.5054. This initially suggested that sparse perturbation features encode mechanism-like structure. When the split was changed to protein-level holdout, the SAE-only macro-AUC dropped to 0.5161 and SAE+LLR to 0.5164 (Table~\ref{tab:mechanism}). This indicates that the variant-level result leaks protein-specific signatures and that current SAE features do not yet support robust cross-protein LOF/GOF/DN prediction.

\begin{table}[h]
\caption{Mechanism classification performance for LOF, GOF and DN labels.}\label{tab:mechanism}
\begin{tabular}{@{}lll@{}}
\toprule
Method & Variant-level macro-AUC & Protein-level macro-AUC \\
\midrule
SAE-LR & 0.7471 & 0.5161 \\
ESM-2 LLR & 0.5054 & 0.4642 \\
SAE+LLR & 0.7488 & 0.5164 \\
\botrule
\end{tabular}
\end{table}

The failed protein-level generalization does not make the perturbation signatures useless. It does mean that the manuscript cannot claim a validated mechanism classifier. The defensible statement is that SAE features contain variant-level mechanism-like structure and produce interpretable hypotheses that require protein-aware validation.

\subsection{ProteinGym shows that SAE disruption is not generic DMS fitness}

We evaluated 217 ProteinGym substitution assays, with finite SAE scores for 214 assays. LLR achieved mean Spearman rho 0.4341. The raw SAE disruption score was negatively correlated with DMS fitness, with mean rho -0.2314, as expected for a damage magnitude. A sign-corrected SAE+LLR ensemble reached mean rho 0.4047, but remained below LLR and beat LLR on only 33.6\% of usable assays (Table~\ref{tab:proteingym}). The classifier-based ProteinGym ensemble is not yet available in the diagnostics.

\begin{table}[h]
\caption{ProteinGym diagnostic results across usable substitution assays.}\label{tab:proteingym}
\begin{tabular}{@{}llll@{}}
\toprule
Score & Assays & Mean Spearman rho & 95\% CI \\
\midrule
ESM-2 LLR & 214 & 0.4341 & [0.4091, 0.4594] \\
SAE disruption & 214 & -0.2314 & [-0.2588, -0.2035] \\
Negated SAE disruption & 214 & 0.2314 & [0.2034, 0.2592] \\
Legacy SAE+LLR & 214 & 0.1993 & [0.1737, 0.2244] \\
Sign-corrected SAE+LLR & 214 & 0.4047 & [0.3818, 0.4282] \\
\botrule
\end{tabular}
\end{table}

This negative result is informative. The SAE perturbation magnitude appears to track disruption-like signal rather than the assay-specific fitness landscapes measured by DMS. Thus ProteinGym should be presented as a diagnostic showing the boundary of the method, not as evidence that SAE+LLR improves generic zero-shot fitness prediction.

\subsection{Indel transfer extends the diagnostic beyond missense substitutions}

We staged 185,655 ClinVar indel records and scored 6,649 reconstructable pathogenic or benign records spanning deletion, insertion, delins and duplication classes. The full run completed for all 6,649 supported records, with 5,274 pathogenic and 1,375 benign labels. The transferred mechanism classifier predicted 4,161 LOF, 1,354 GOF and 1,134 DN records. The SAE damage-score AUC for indel pathogenicity was 0.7735.

The current indel result is not a final benchmark because the score is based on an affected-region SAE-delta heuristic and remains below the internal AUC target of 0.85. Nevertheless, it is a useful direction because scalar missense tools do not uniformly cover insertions, deletions and duplications. The near-term experimental value is to use indel perturbation signatures to nominate variants for targeted assays in genes with tractable readouts.

\subsection{Circuit-latent dictionaries expose generative features but do not yet support steering claims}

For the generative branch, we retrained ProtGPT2 and ZymCTRL circuit-latent dictionaries with $d_\mathrm{CLT}=8192$, $k=128$ and 200,000 training steps. ProtGPT2 v2 reached mean FVU 0.2013 with alive fraction 0.3803. ZymCTRL v2 reached mean FVU 0.3307 with alive fraction 0.3318. These checkpoints improved the earlier weak CLTs but retained high dead-feature fractions.

ZymCTRL steering was tested across eight enzyme classes with 200 steered and 200 unsteered generations per condition, using layers 3, 12 and 30, five features per layer and multiplier 2.5. No enzyme class showed a significant positive steering effect. A stronger multiplier 5.0 sweep over lipase, kinase and carbonic anhydrase also produced no significant positive class; lipase showed a numerical shift of +0.0717 but permutation $P=0.1066$.

The lysozyme case study generated 200 steered sequences and retained ten leads. The simple Glu/Asp motif filter was positive in 97.5\% of generated records, but structural quality did not favor steering: ESMFold mean pLDDT was 67.7 for ten steered leads and 73.2 for 20 unsteered baseline sequences, with globally confident fractions of 0.60 and 0.70, respectively. A causal ablation test on ZymCTRL layer 12 feature 8088 gave exactly zero mean absolute log-probability change across 1,665 positions, confirming that ablation claims require stricter hook diagnostics. Cross-model conservation between ZymCTRL and ProGen2-medium showed nontrivial feature-space similarity at mapped layers, but this does not rescue the steering claim.

\section{Discussion}\label{sec:discussion}

The current evidence supports BioCC as an interpretable perturbation and circuit diagnostic, not as a finished predictive replacement for specialized variant-effect models. For R1, sparse features add measurable information to ESM-2 LLR and provide residue-linked feature tables that scalar tools do not expose. This is enough to justify follow-up as a wet-lab prioritization framework. The strongest immediate wet-lab use is to select variants with similar scalar pathogenicity but divergent SAE perturbation profiles, then test whether these profiles correspond to different biochemical mechanisms.

The evidence does not yet support stronger claims. AlphaMissense is substantially better on scalar pathogenicity in the current cohorts. Mechanism prediction fails under protein-level holdout. ProteinGym does not show a mean advantage over LLR. The indel result is promising in scope but preliminary in accuracy. These limitations should be kept in the manuscript because they clarify the biological question that remains: whether SAE perturbation signatures enrich experimentally separable mechanism classes even when they do not function as a generic classifier.

For R2, the circuit-latent workflow is useful for diagnostics and layer mapping, but the current steering experiments are negative. A Nature Methods-level paper would need either a calibrated steering metric stack with significant causal effects or a reframed contribution centered on interpretable model auditing. Until then, R2 should be treated as a secondary demonstration of the same diagnostic philosophy rather than as a validated protein-design pipeline.

The next manuscript-critical experiments are clear. First, complete the R1 head-to-head table with ESM-1v, gMVP, dbNSFP-derived predictors and PrimateAI-3D where licenses allow. Second, perform residual analysis against AlphaMissense and LLR to identify variants where SAE features add orthogonal information. Third, curate a small wet-lab-ready variant panel in proteins with direct readouts, such as ion channels, enzymes or cancer signaling proteins. Fourth, either retire the R2 steering claim or rerun it only after hook diagnostics and external EC metrics are calibrated.

\section{Online Methods}\label{sec:methods}

\subsection{ESM-2 sparse autoencoder training}

SAEs were trained on ESM-2-3B hidden activations from layers 19, 23, 27, 31 and 35. The production checkpoints used sparse dictionaries with 16,384 features and TopK activation. Training used UniRef50 protein sequences and frozen ESM-2 activations. Early experiments identified dead-feature spirals under BatchTopK and a pre-TopK ReLU issue; the production analysis therefore uses the corrected SAE checkpoints documented in the R1 experiment log.

\subsection{Feature annotation}

For each SAE layer, features were annotated by comparing top firing residue positions with Swiss-Prot features, Pfam families, Gene Ontology labels and BioLiP/PDB binding-site proxies. Feature-label alignment was summarized by F1 score. The firing-position rerun stores top residue-level examples for manual inspection and enables expanded annotation tables. Known, partial and useful feature counts were defined by F1 thresholds of 0.5, 0.2 and 0.1, respectively.

\subsection{Variant perturbation scoring}

For a protein sequence and missense variant, ESM-2 activations were computed for the wild-type and mutant sequences. Each layer-specific SAE encoded both activation tensors. The variant perturbation signature consisted of feature activation deltas across the five production layers. We summarized deltas as raw disruption magnitudes, annotation-selected feature vectors and learned logistic-regression scores. LLR was computed as the ESM-2 masked-marginal log-likelihood ratio between mutant and wild-type residues.

\subsection{Pathogenicity cohorts and baselines}

The primary local cohort contained 2,000 ClinVar missense variants balanced between pathogenic and benign labels. A 101-variant cancer holdout was used as an additional generalization set. Models were evaluated by ROC AUC with bootstrap 95\% confidence intervals. Local baselines included ESM-2 LLR, SAE-LR and SAE+LLR z-sum. AlphaMissense was matched by UniProt identifier and protein substitution. At the time of this draft, PrimateAI-3D remained gated, while ESM-1v, gMVP and dbNSFP resources were being staged for a unified head-to-head table.

\subsection{Mechanism label construction and classification}

LOF, GOF and DN labels were assembled from curated mechanism tables and gene-level fallbacks where available. The resulting labeled subset contained 497 variants across 255 proteins. Multinomial logistic regression was trained on annotation-selected SAE features, LLR alone or concatenated SAE+LLR features. Variant-level cross-validation split individual variants, whereas protein-level holdout grouped variants by protein to prevent leakage of protein-specific signatures. The protein-level result is treated as the decision metric.

\subsection{ProteinGym benchmark}

ProteinGym substitution assays were scored with ESM-2 LLR and SAE perturbation scores. The primary metric was per-assay Spearman correlation between model score and experimental DMS score, followed by the mean over usable assays and bootstrap confidence intervals. Because SAE disruption magnitude is expected to increase as fitness decreases, we evaluated both the raw sign and a sign-corrected ensemble.

\subsection{Indel transfer scoring}

ClinVar indels were parsed into reconstructable protein-sequence edits where transcript context was sufficient for deletion, insertion, duplication or delins reconstruction. For each supported record, an affected-region SAE perturbation signature was computed and passed through the missense-trained mechanism classifier or damage-score heuristic. Frameshift records without reliable protein reconstruction were excluded from the current analysis.

\subsection{Circuit-latent training and generative diagnostics}

Circuit-latent dictionaries were trained for ProtGPT2 and ZymCTRL using hidden activations from generative PLMs. The v2 checkpoints used 36 layers, $d_\mathrm{CLT}=8192$, $k=128$ and 200,000 training steps. Quick evaluation measured fraction of variance unexplained, L0 and dead/alive feature fractions on held-out sequences. ZymCTRL enzyme-class steering selected top direct-effect features for eight EC classes and applied activation multipliers during generation. Steering effects were measured by class-purity scores with bootstrap confidence intervals and permutation tests. ESMFold was used for preliminary structural QC of lysozyme lead sequences.

\subsection{Statistics}

ROC AUCs were reported with nonparametric bootstrap 95\% confidence intervals. ProteinGym correlations used Spearman rho per assay and bootstrap confidence intervals over assays. Steering benchmarks used permutation tests over steered and unsteered generations. Where protein-level and variant-level cross-validation disagreed, the protein-level split was used for the scientific interpretation.

\section*{Data availability}

All datasets used in this draft are public or license-gated public resources staged locally for analysis. Exact local paths are recorded in the repository status files and experiment logs. Large checkpoints and external databases remain on H200/OSS storage. A submission version should replace this paragraph with stable accession numbers, download URLs and license statements.

\section*{Code availability}

The current analysis code is in the BioCC repository under \texttt{Research1/} and \texttt{Research2/}. A submission version should provide a clean tagged release, frozen environment file, command log and minimal reproducible examples for variant scoring and circuit diagnostics.

\section*{Acknowledgements}

TODO: add funding, compute and institutional acknowledgements. H200 experiments were run on the remote Hangzhou cluster; local analyses were staged in the BioCC repository.

\section*{Author contributions}

TODO: update after coauthor list is finalized.

\section*{Competing interests}

The authors declare no competing interests. TODO: confirm before submission.

\bibliography{references}

\end{document}
